% Apéndice: Manual de usuario — Sistema de Seguimiento de Egresados
% Incluir en el documento principal con: % Apéndice: Manual de usuario — Sistema de Seguimiento de Egresados
% Incluir en el documento principal con: % Apéndice: Manual de usuario — Sistema de Seguimiento de Egresados
% Incluir en el documento principal con: % Apéndice: Manual de usuario — Sistema de Seguimiento de Egresados
% Incluir en el documento principal con: \input{docs/manual_usuario.tex}
% Requiere en el preámbulo: \usepackage{float} \usepackage{graphicx} \usepackage{hyperref}

\chapter{Manual de usuarios}
\label{ap:manual-usuario}

\begin{center}
\textbf{MANUAL DE USO DEL SISTEMA DE SEGUIMIENTO DE EGRESADOS}\\[0.5em]
\textit{Trabajo de Título --- SeguimientoApp}
\end{center}

\section{Introducción}
\label{sec:manual-intro}

El objetivo de realizar este manual es poder ayudar a los diferentes usuarios que quieran utilizar este sistema, guiándolos en los pasos que debe realizar ya sea para ingresar al sistema como en el uso del mismo.

Este proyecto fue desarrollado con el fin de dar seguimiento a la trayectoria profesional de los egresados de la institución, consolidar información laboral y tecnológica, y apoyar la toma de decisiones académicas e institucionales mediante reportes y encuestas. El sistema contempla dos perfiles principales: \textbf{egresado} y \textbf{administrador}.

A continuación se detalla, en la primera sección, el ingreso al sistema; en las secciones siguientes, el funcionamiento de cada módulo según el rol del usuario.

\subsection{Ingreso al sistema}
\label{subsec:manual-ingreso}

\subsubsection{Requisitos y acceso general}

Para utilizar el sistema, el usuario debe abrir un navegador web y dirigirse a la URL de despliegue de la aplicación (por ejemplo, la ruta \texttt{/auth} del frontend). En la pantalla de inicio se podrá alternar entre el \textbf{Portal Egresados} y el \textbf{Panel Administrativo} mediante el botón de cambio de rol, como se muestra en la Figura~\ref{fig:login}.

\begin{figure}[H]
\centering
% \includegraphics[width=0.85\textwidth]{figuras/manual/login.png}
\fbox{\parbox{0.85\textwidth}{\centering\vspace{2cm}[Insertar captura: pantalla de login]\vspace{2cm}}}
\caption{Pantalla de inicio de sesión con selector de rol.}
\label{fig:login}
\end{figure}

\subsubsection{Ingreso del usuario egresado}

El acceso del egresado al sistema se realiza mediante autenticación con LinkedIn y contempla, para un usuario nuevo, un flujo de cuatro etapas que se describen a continuación: ingreso a la aplicación, autorización en LinkedIn, completitud del perfil profesional y acceso al dashboard principal.

\paragraph{Paso 1: Ingreso a la página del sistema.}

El primer paso consiste en abrir un navegador web y dirigirse a la URL de despliegue de la aplicación (por ejemplo, \texttt{http://localhost:4200} en ambiente de desarrollo, o la dirección institucional definida en producción). Al cargar el sitio, el sistema redirige automáticamente a la ruta de autenticación \texttt{/auth}, donde se presenta la pantalla denominada \emph{Portal Egresados}, como se muestra en la Figura~\ref{fig:login-egresado-paso1}.

En esta pantalla, el usuario debe verificar que \textbf{no} esté activa la vista de \emph{Panel Administrativo}. Si el formulario muestra campos de correo y contraseña, debe presionar el botón superior \emph{Egresados} para volver al modo correspondiente al perfil de egresado. Los elementos visibles en esta etapa son los siguientes:

\begin{enumerate}
\item \textbf{Encabezado del portal:} título \emph{Portal Egresados} y mensaje de bienvenida.
\item \textbf{Ícono de LinkedIn} junto al texto \emph{Conecta con tu perfil profesional}, que indica el método de acceso habilitado para este rol.
\item \textbf{Botón Ingresar con LinkedIn:} acción principal que inicia el proceso de autenticación externa.
\item \textbf{Botón de cambio de rol} (parte superior): permite alternar entre egresado y administrador; debe permanecer en modo egresado.
\end{enumerate}

\begin{figure}[H]
\centering
\fbox{\parbox{0.85\textwidth}{\centering\vspace{2cm}[Insertar captura: Portal Egresados --- Paso 1]\vspace{2cm}}}
\caption{Paso 1 --- Página de ingreso del egresado en la URL del sistema.}
\label{fig:login-egresado-paso1}
\end{figure}

\paragraph{Paso 2: Redirección y autenticación en LinkedIn.}

Al presionar el botón \emph{Ingresar con LinkedIn}, el sistema solicita al servidor la URL de autorización de LinkedIn y muestra un cuadro de diálogo con el mensaje \emph{Conectando con LinkedIn} o \emph{Redirigiendo a LinkedIn}, indicando que la operación está en curso. Posteriormente, el navegador abandona temporalmente la aplicación y carga el sitio oficial de LinkedIn (\texttt{linkedin.com}), donde el usuario deberá:

\begin{enumerate}
\item Iniciar sesión con su cuenta de LinkedIn, si aún no lo ha hecho en esa sesión del navegador.
\item Revisar los permisos que la aplicación solicita (acceso al perfil básico y datos profesionales).
\item Presionar el botón \emph{Permitir} o \emph{Aceptar} para autorizar la vinculación de su cuenta con el sistema.
\end{enumerate}

Una vez autorizado el acceso, LinkedIn redirige nuevamente a la aplicación con un código de autenticación. El sistema procesa dicho código en segundo plano, registra al usuario si es su primer ingreso o actualiza sus datos si ya existía previamente, y muestra un mensaje de éxito (\emph{Accediendo al portal...}). Finalmente, el egresado es conducido al área privada del sistema bajo la ruta \texttt{/user}, como se aprecia en la Figura~\ref{fig:login-egresado-paso2}.

\begin{figure}[H]
\centering
\fbox{\parbox{0.85\textwidth}{\centering\vspace{2cm}[Insertar captura: autorización LinkedIn --- Paso 2]\vspace{2cm}}}
\caption{Paso 2 --- Pantalla de autorización en LinkedIn y retorno al sistema.}
\label{fig:login-egresado-paso2}
\end{figure}

Si el usuario cancela la autorización o ocurre un error en la autenticación, el sistema permanece en la pantalla de login y despliega un mensaje de error descriptivo, permitiendo reintentar el proceso.

\paragraph{Paso 3: Completar el perfil profesional (usuario nuevo).}

Cuando el egresado ingresa por primera vez, el sistema importa desde LinkedIn información básica como nombre, fotografía, correo, empresa, cargo, ubicación y resumen profesional. Sin embargo, para generar métricas precisas del mercado laboral, la aplicación exige completar un conjunto de \textbf{campos obligatorios} antes de habilitar el dashboard y el módulo de encuestas.

Si el perfil no está completo, al intentar acceder a \texttt{/user/dashboard} o a las encuestas, el sistema redirige automáticamente a la ruta \texttt{/user/mi-perfil} y muestra un aviso destacado con el mensaje \emph{Complete su perfil profesional}, junto a una barra de progreso de completitud (ver Figura~\ref{fig:login-egresado-paso3}). Los campos que el usuario debe completar son:

\begin{enumerate}
\item Años de experiencia laboral.
\item Nivel de educación.
\item Especialidad técnica.
\item Tipo de empleo actual.
\item Disponibilidad para cambio de trabajo.
\item Área de interés.
\end{enumerate}

Además, el formulario incluye secciones de \emph{Información personal} (nombre, ubicación), \emph{Información profesional} (posición, empresa, industria, resumen) y \emph{Datos para métricas} (tecnologías, rango salarial, satisfacción laboral, entre otros campos opcionales). El usuario debe completar los campos marcados con asterisco rojo ($*$), presionar el botón \emph{Guardar cambios} y verificar que el indicador de progreso alcance el \textbf{100\%}, momento en el cual aparece el mensaje \emph{¡Perfil completado exitosamente!}

\begin{figure}[H]
\centering
\fbox{\parbox{0.85\textwidth}{\centering\vspace{2cm}[Insertar captura: completar perfil --- Paso 3]\vspace{2cm}}}
\caption{Paso 3 --- Completitud del perfil profesional para usuario nuevo.}
\label{fig:login-egresado-paso3}
\end{figure}

Hasta que el perfil no se encuentre completo, el menú lateral seguirá visible, pero las rutas protegidas del dashboard y encuestas permanecerán bloqueadas por el sistema de validación de perfil.

\paragraph{Paso 4: Acceso al Dashboard principal.}

Una vez completado el perfil, el egresado puede acceder al \emph{Dashboard principal} seleccionando en el menú lateral la opción \emph{Reporte egresados} $\rightarrow$ \emph{Dashboard principal}, o navegando directamente a \texttt{/user/dashboard}. En esta pantalla el usuario visualiza por primera vez el conjunto de indicadores agregados del mercado laboral tecnológico, alimentados con los datos consolidados de la comunidad de egresados (ver Figura~\ref{fig:login-egresado-paso4}).

Los elementos principales del dashboard son:

\begin{enumerate}
\item \textbf{Encabezado personalizado:} muestra la fotografía importada de LinkedIn, el nombre del egresado, su cargo actual, ubicación, empresa e industria. Al hacer clic en el nombre o la imagen, el usuario puede volver a \emph{Mi perfil}.
\item \textbf{Tarjetas de resumen del mercado:} indicadores como profesionales activos, satisfacción laboral promedio, tecnologías más utilizadas y empleabilidad del sector.
\item \textbf{Gráficos interactivos:} distribución salarial, tecnologías demandadas, análisis de experiencia, mapa de calor por industria y métricas avanzadas, accesibles también desde los submenús del panel lateral.
\item \textbf{Botón de menú (\texttt{☰}):} permite desplegar el panel de navegación para acceder a encuestas u otras secciones del reporte.
\item \textbf{Botón volver arriba:} aparece al desplazarse hacia abajo en páginas extensas de gráficos.
\end{enumerate}

\begin{figure}[H]
\centering
\fbox{\parbox{0.85\textwidth}{\centering\vspace{2cm}[Insertar captura: Dashboard principal --- Paso 4]\vspace{2cm}}}
\caption{Paso 4 --- Dashboard principal del egresado tras completar el perfil.}
\label{fig:login-egresado-paso4}
\end{figure}

A partir de este punto, el egresado dispone de acceso completo a las funcionalidades de su rol: consulta de reportes, respuesta de encuestas y actualización permanente de su perfil profesional (detalladas en la Sección~\ref{sec:manual-egresado}).

\subsubsection{Ingreso del usuario administrador}

El administrador debe activar la vista \emph{Panel Administrativo}, ingresar correo electrónico y contraseña, y presionar \emph{Iniciar Sesión}. La recuperación de contraseña está disponible mediante el enlace \emph{¿Olvidaste tu contraseña?}, que despliega un asistente en tres pasos: solicitud de código, verificación e ingreso de nueva contraseña.

\begin{figure}[H]
\centering
\fbox{\parbox{0.85\textwidth}{\centering\vspace{2cm}[Insertar captura: login administrador]\vspace{2cm}}}
\caption{Formulario de acceso administrativo.}
\label{fig:login-admin}
\end{figure}

\subsubsection{Navegación general}

Una vez autenticado, la interfaz se compone de una barra superior, un botón de menú (\texttt{☰}) que despliega el panel lateral, y el área de contenido. En el menú lateral, la tarjeta de perfil permite acceder a \emph{Mi perfil}. Las opciones del menú dependen del rol del usuario.

\begin{figure}[H]
\centering
\fbox{\parbox{0.85\textwidth}{\centering\vspace{2cm}[Insertar captura: layout principal]\vspace{2cm}}}
\caption{Vista general del sistema con menú lateral.}
\label{fig:layout}
\end{figure}

\section{Módulos del sistema --- Usuario administrador}
\label{sec:manual-admin}

\subsection{Dashboard y reporte de egresados}
\label{subsec:dashboard-admin}

Para ingresar al dashboard, el administrador debe seleccionar en el menú lateral la opción \emph{Reporte egresados} $\rightarrow$ \emph{Dashboard principal} (ruta \texttt{/admin/dashboard}).

La página presenta métricas agregadas del mercado laboral tecnológico. Los elementos principales son:

\begin{enumerate}
\item Encabezado con indicadores generales del mercado laboral.
\item Tarjeta de \emph{Nuevos registros} con la cantidad de egresados incorporados recientemente.
\item Conjunto de gráficos interactivos alimentados por los datos de los perfiles.
\end{enumerate}

Desde el submenú \emph{Reporte egresados} es posible acceder directamente a secciones como análisis salarial, tecnologías más demandadas, satisfacción laboral, mapa de calor y métricas avanzadas; el sistema desplaza la vista hasta el gráfico seleccionado.

\begin{figure}[H]
\centering
\fbox{\parbox{0.85\textwidth}{\centering\vspace{2cm}[Insertar captura: dashboard administrador]\vspace{2cm}}}
\caption{Dashboard principal del administrador.}
\label{fig:dashboard-admin}
\end{figure}

\subsection{Gestión de usuarios egresados}

\subsubsection{Usuarios activos}

Ruta: \texttt{/admin/view-users}. La tabla lista los egresados registrados. Los elementos de la interfaz son:

\begin{enumerate}
\item \textbf{Buscador:} permite filtrar por nombre, correo o empresa.
\item \textbf{Botón limpiar filtros:} restablece la búsqueda.
\item \textbf{Botón actualizar:} recarga los datos desde el servidor.
\item \textbf{Tabla paginada:} muestra el detalle de cada egresado y acciones de edición o desactivación.
\end{enumerate}

\begin{figure}[H]
\centering
\fbox{\parbox{0.85\textwidth}{\centering\vspace{2cm}[Insertar captura: usuarios activos]\vspace{2cm}}}
\caption{Listado de usuarios egresados activos.}
\label{fig:usuarios-activos}
\end{figure}

\subsubsection{Usuarios eliminados}

Ruta: \texttt{/admin/view-deleted-users}. Muestra egresados desactivados y permite su reactivación.

\subsection{Gestión de administradores}

Rutas: \texttt{/admin/view-admin} (activos) y \texttt{/admin/view-deleted-admin} (eliminados). Permite administrar cuentas con privilegios administrativos de forma análoga a la gestión de egresados.

\subsection{Gestión de encuestas}

\subsubsection{Crear o editar encuesta}

Ruta: \texttt{/admin/create-encuesta}. El formulario se organiza en: (1) datos principales (título, estado, descripción, fechas); (2) definición de preguntas y tipos de respuesta; (3) guardado. Para editar una encuesta existente se utiliza \texttt{/admin/editar-encuesta/:id}.

\subsubsection{Listado y resultados}

\begin{itemize}
\item \textbf{Mis encuestas} (\texttt{/admin/view-encuesta}): listado y mantenimiento de encuestas.
\item \textbf{Resultados} (\texttt{/admin/view-encuestas-resultados}): análisis de respuestas, gráficos y exportación de datos.
\end{itemize}

\begin{figure}[H]
\centering
\fbox{\parbox{0.85\textwidth}{\centering\vspace{2cm}[Insertar captura: resultados encuesta]\vspace{2cm}}}
\caption{Vista de resultados de encuestas.}
\label{fig:resultados-encuesta}
\end{figure}

\subsection{Mi perfil del administrador}

Ruta: \texttt{/admin/mi-perfil}. Permite actualizar datos personales y credenciales del administrador.

\section{Módulos del sistema --- Usuario egresado}
\label{sec:manual-egresado}

\subsection{Completar perfil profesional}
\label{subsec:perfil-egresado}

Tras el primer ingreso, el sistema puede restringir el acceso al dashboard y a encuestas hasta completar el perfil (ruta \texttt{/user/mi-perfil}). Los campos obligatorios son: años de experiencia, nivel de educación, especialidad técnica, tipo de empleo actual, disponibilidad de cambio y área de interés.

\begin{figure}[H]
\centering
\fbox{\parbox{0.85\textwidth}{\centering\vspace{2cm}[Insertar captura: perfil incompleto]\vspace{2cm}}}
\caption{Solicitud de completar perfil profesional.}
\label{fig:perfil-incompleto}
\end{figure}

\subsection{Dashboard del egresado}

Ruta: \texttt{/user/dashboard}. Ofrece las mismas visualizaciones agregadas que el panel administrativo, sin la tarjeta de nuevos registros. El egresado puede navegar por las subsecciones del menú \emph{Reporte egresados}.

\subsection{Encuestas}

\subsubsection{Mis encuestas}

Ruta: \texttt{/user/view-encuestas}. Lista encuestas disponibles con buscador y filtro por estado. El botón \emph{Responder} conduce a \texttt{/user/responder-encuesta/:id}.

\subsubsection{Encuestas completadas}

Ruta: \texttt{/user/encuesta-completada}. Historial de encuestas ya enviadas.

\begin{figure}[H]
\centering
\fbox{\parbox{0.85\textwidth}{\centering\vspace{2cm}[Insertar captura: responder encuesta]\vspace{2cm}}}
\caption{Formulario de respuesta de encuesta.}
\label{fig:responder-encuesta}
\end{figure}

\subsection{Mi perfil del egresado}

Ruta: \texttt{/user/mi-perfil}. Permite editar y guardar la información profesional que alimenta los reportes del sistema.

\section{Resumen de rutas}

La Tabla~\ref{tab:rutas-manual} resume las rutas principales por rol.

\begin{table}[H]
\centering
\caption{Rutas principales del sistema SeguimientoApp.}
\label{tab:rutas-manual}
\begin{tabular}{|l|l|l|}
\hline
\textbf{Funcionalidad} & \textbf{Administrador} & \textbf{Egresado} \\
\hline
Inicio de sesión & \texttt{/auth} & \texttt{/auth} (LinkedIn) \\
Dashboard & \texttt{/admin/dashboard} & \texttt{/user/dashboard} \\
Mi perfil & \texttt{/admin/mi-perfil} & \texttt{/user/mi-perfil} \\
Usuarios & \texttt{/admin/view-users} & --- \\
Encuestas & \texttt{/admin/view-encuesta} & \texttt{/user/view-encuestas} \\
Resultados & \texttt{/admin/view-encuestas-resultados} & --- \\
Responder encuesta & --- & \texttt{/user/responder-encuesta/:id} \\
\hline
\end{tabular}
\end{table}

% Requiere en el preámbulo: \usepackage{float} \usepackage{graphicx} \usepackage{hyperref}

\chapter{Manual de usuarios}
\label{ap:manual-usuario}

\begin{center}
\textbf{MANUAL DE USO DEL SISTEMA DE SEGUIMIENTO DE EGRESADOS}\\[0.5em]
\textit{Trabajo de Título --- SeguimientoApp}
\end{center}

\section{Introducción}
\label{sec:manual-intro}

El objetivo de realizar este manual es poder ayudar a los diferentes usuarios que quieran utilizar este sistema, guiándolos en los pasos que debe realizar ya sea para ingresar al sistema como en el uso del mismo.

Este proyecto fue desarrollado con el fin de dar seguimiento a la trayectoria profesional de los egresados de la institución, consolidar información laboral y tecnológica, y apoyar la toma de decisiones académicas e institucionales mediante reportes y encuestas. El sistema contempla dos perfiles principales: \textbf{egresado} y \textbf{administrador}.

A continuación se detalla, en la primera sección, el ingreso al sistema; en las secciones siguientes, el funcionamiento de cada módulo según el rol del usuario.

\subsection{Ingreso al sistema}
\label{subsec:manual-ingreso}

\subsubsection{Requisitos y acceso general}

Para utilizar el sistema, el usuario debe abrir un navegador web y dirigirse a la URL de despliegue de la aplicación (por ejemplo, la ruta \texttt{/auth} del frontend). En la pantalla de inicio se podrá alternar entre el \textbf{Portal Egresados} y el \textbf{Panel Administrativo} mediante el botón de cambio de rol, como se muestra en la Figura~\ref{fig:login}.

\begin{figure}[H]
\centering
% \includegraphics[width=0.85\textwidth]{figuras/manual/login.png}
\fbox{\parbox{0.85\textwidth}{\centering\vspace{2cm}[Insertar captura: pantalla de login]\vspace{2cm}}}
\caption{Pantalla de inicio de sesión con selector de rol.}
\label{fig:login}
\end{figure}

\subsubsection{Ingreso del usuario egresado}

El acceso del egresado al sistema se realiza mediante autenticación con LinkedIn y contempla, para un usuario nuevo, un flujo de cuatro etapas que se describen a continuación: ingreso a la aplicación, autorización en LinkedIn, completitud del perfil profesional y acceso al dashboard principal.

\paragraph{Paso 1: Ingreso a la página del sistema.}

El primer paso consiste en abrir un navegador web y dirigirse a la URL de despliegue de la aplicación (por ejemplo, \texttt{http://localhost:4200} en ambiente de desarrollo, o la dirección institucional definida en producción). Al cargar el sitio, el sistema redirige automáticamente a la ruta de autenticación \texttt{/auth}, donde se presenta la pantalla denominada \emph{Portal Egresados}, como se muestra en la Figura~\ref{fig:login-egresado-paso1}.

En esta pantalla, el usuario debe verificar que \textbf{no} esté activa la vista de \emph{Panel Administrativo}. Si el formulario muestra campos de correo y contraseña, debe presionar el botón superior \emph{Egresados} para volver al modo correspondiente al perfil de egresado. Los elementos visibles en esta etapa son los siguientes:

\begin{enumerate}
\item \textbf{Encabezado del portal:} título \emph{Portal Egresados} y mensaje de bienvenida.
\item \textbf{Ícono de LinkedIn} junto al texto \emph{Conecta con tu perfil profesional}, que indica el método de acceso habilitado para este rol.
\item \textbf{Botón Ingresar con LinkedIn:} acción principal que inicia el proceso de autenticación externa.
\item \textbf{Botón de cambio de rol} (parte superior): permite alternar entre egresado y administrador; debe permanecer en modo egresado.
\end{enumerate}

\begin{figure}[H]
\centering
\fbox{\parbox{0.85\textwidth}{\centering\vspace{2cm}[Insertar captura: Portal Egresados --- Paso 1]\vspace{2cm}}}
\caption{Paso 1 --- Página de ingreso del egresado en la URL del sistema.}
\label{fig:login-egresado-paso1}
\end{figure}

\paragraph{Paso 2: Redirección y autenticación en LinkedIn.}

Al presionar el botón \emph{Ingresar con LinkedIn}, el sistema solicita al servidor la URL de autorización de LinkedIn y muestra un cuadro de diálogo con el mensaje \emph{Conectando con LinkedIn} o \emph{Redirigiendo a LinkedIn}, indicando que la operación está en curso. Posteriormente, el navegador abandona temporalmente la aplicación y carga el sitio oficial de LinkedIn (\texttt{linkedin.com}), donde el usuario deberá:

\begin{enumerate}
\item Iniciar sesión con su cuenta de LinkedIn, si aún no lo ha hecho en esa sesión del navegador.
\item Revisar los permisos que la aplicación solicita (acceso al perfil básico y datos profesionales).
\item Presionar el botón \emph{Permitir} o \emph{Aceptar} para autorizar la vinculación de su cuenta con el sistema.
\end{enumerate}

Una vez autorizado el acceso, LinkedIn redirige nuevamente a la aplicación con un código de autenticación. El sistema procesa dicho código en segundo plano, registra al usuario si es su primer ingreso o actualiza sus datos si ya existía previamente, y muestra un mensaje de éxito (\emph{Accediendo al portal...}). Finalmente, el egresado es conducido al área privada del sistema bajo la ruta \texttt{/user}, como se aprecia en la Figura~\ref{fig:login-egresado-paso2}.

\begin{figure}[H]
\centering
\fbox{\parbox{0.85\textwidth}{\centering\vspace{2cm}[Insertar captura: autorización LinkedIn --- Paso 2]\vspace{2cm}}}
\caption{Paso 2 --- Pantalla de autorización en LinkedIn y retorno al sistema.}
\label{fig:login-egresado-paso2}
\end{figure}

Si el usuario cancela la autorización o ocurre un error en la autenticación, el sistema permanece en la pantalla de login y despliega un mensaje de error descriptivo, permitiendo reintentar el proceso.

\paragraph{Paso 3: Completar el perfil profesional (usuario nuevo).}

Cuando el egresado ingresa por primera vez, el sistema importa desde LinkedIn información básica como nombre, fotografía, correo, empresa, cargo, ubicación y resumen profesional. Sin embargo, para generar métricas precisas del mercado laboral, la aplicación exige completar un conjunto de \textbf{campos obligatorios} antes de habilitar el dashboard y el módulo de encuestas.

Si el perfil no está completo, al intentar acceder a \texttt{/user/dashboard} o a las encuestas, el sistema redirige automáticamente a la ruta \texttt{/user/mi-perfil} y muestra un aviso destacado con el mensaje \emph{Complete su perfil profesional}, junto a una barra de progreso de completitud (ver Figura~\ref{fig:login-egresado-paso3}). Los campos que el usuario debe completar son:

\begin{enumerate}
\item Años de experiencia laboral.
\item Nivel de educación.
\item Especialidad técnica.
\item Tipo de empleo actual.
\item Disponibilidad para cambio de trabajo.
\item Área de interés.
\end{enumerate}

Además, el formulario incluye secciones de \emph{Información personal} (nombre, ubicación), \emph{Información profesional} (posición, empresa, industria, resumen) y \emph{Datos para métricas} (tecnologías, rango salarial, satisfacción laboral, entre otros campos opcionales). El usuario debe completar los campos marcados con asterisco rojo ($*$), presionar el botón \emph{Guardar cambios} y verificar que el indicador de progreso alcance el \textbf{100\%}, momento en el cual aparece el mensaje \emph{¡Perfil completado exitosamente!}

\begin{figure}[H]
\centering
\fbox{\parbox{0.85\textwidth}{\centering\vspace{2cm}[Insertar captura: completar perfil --- Paso 3]\vspace{2cm}}}
\caption{Paso 3 --- Completitud del perfil profesional para usuario nuevo.}
\label{fig:login-egresado-paso3}
\end{figure}

Hasta que el perfil no se encuentre completo, el menú lateral seguirá visible, pero las rutas protegidas del dashboard y encuestas permanecerán bloqueadas por el sistema de validación de perfil.

\paragraph{Paso 4: Acceso al Dashboard principal.}

Una vez completado el perfil, el egresado puede acceder al \emph{Dashboard principal} seleccionando en el menú lateral la opción \emph{Reporte egresados} $\rightarrow$ \emph{Dashboard principal}, o navegando directamente a \texttt{/user/dashboard}. En esta pantalla el usuario visualiza por primera vez el conjunto de indicadores agregados del mercado laboral tecnológico, alimentados con los datos consolidados de la comunidad de egresados (ver Figura~\ref{fig:login-egresado-paso4}).

Los elementos principales del dashboard son:

\begin{enumerate}
\item \textbf{Encabezado personalizado:} muestra la fotografía importada de LinkedIn, el nombre del egresado, su cargo actual, ubicación, empresa e industria. Al hacer clic en el nombre o la imagen, el usuario puede volver a \emph{Mi perfil}.
\item \textbf{Tarjetas de resumen del mercado:} indicadores como profesionales activos, satisfacción laboral promedio, tecnologías más utilizadas y empleabilidad del sector.
\item \textbf{Gráficos interactivos:} distribución salarial, tecnologías demandadas, análisis de experiencia, mapa de calor por industria y métricas avanzadas, accesibles también desde los submenús del panel lateral.
\item \textbf{Botón de menú (\texttt{☰}):} permite desplegar el panel de navegación para acceder a encuestas u otras secciones del reporte.
\item \textbf{Botón volver arriba:} aparece al desplazarse hacia abajo en páginas extensas de gráficos.
\end{enumerate}

\begin{figure}[H]
\centering
\fbox{\parbox{0.85\textwidth}{\centering\vspace{2cm}[Insertar captura: Dashboard principal --- Paso 4]\vspace{2cm}}}
\caption{Paso 4 --- Dashboard principal del egresado tras completar el perfil.}
\label{fig:login-egresado-paso4}
\end{figure}

A partir de este punto, el egresado dispone de acceso completo a las funcionalidades de su rol: consulta de reportes, respuesta de encuestas y actualización permanente de su perfil profesional (detalladas en la Sección~\ref{sec:manual-egresado}).

\subsubsection{Ingreso del usuario administrador}

El administrador debe activar la vista \emph{Panel Administrativo}, ingresar correo electrónico y contraseña, y presionar \emph{Iniciar Sesión}. La recuperación de contraseña está disponible mediante el enlace \emph{¿Olvidaste tu contraseña?}, que despliega un asistente en tres pasos: solicitud de código, verificación e ingreso de nueva contraseña.

\begin{figure}[H]
\centering
\fbox{\parbox{0.85\textwidth}{\centering\vspace{2cm}[Insertar captura: login administrador]\vspace{2cm}}}
\caption{Formulario de acceso administrativo.}
\label{fig:login-admin}
\end{figure}

\subsubsection{Navegación general}

Una vez autenticado, la interfaz se compone de una barra superior, un botón de menú (\texttt{☰}) que despliega el panel lateral, y el área de contenido. En el menú lateral, la tarjeta de perfil permite acceder a \emph{Mi perfil}. Las opciones del menú dependen del rol del usuario.

\begin{figure}[H]
\centering
\fbox{\parbox{0.85\textwidth}{\centering\vspace{2cm}[Insertar captura: layout principal]\vspace{2cm}}}
\caption{Vista general del sistema con menú lateral.}
\label{fig:layout}
\end{figure}

\section{Módulos del sistema --- Usuario administrador}
\label{sec:manual-admin}

\subsection{Dashboard y reporte de egresados}
\label{subsec:dashboard-admin}

Para ingresar al dashboard, el administrador debe seleccionar en el menú lateral la opción \emph{Reporte egresados} $\rightarrow$ \emph{Dashboard principal} (ruta \texttt{/admin/dashboard}).

La página presenta métricas agregadas del mercado laboral tecnológico. Los elementos principales son:

\begin{enumerate}
\item Encabezado con indicadores generales del mercado laboral.
\item Tarjeta de \emph{Nuevos registros} con la cantidad de egresados incorporados recientemente.
\item Conjunto de gráficos interactivos alimentados por los datos de los perfiles.
\end{enumerate}

Desde el submenú \emph{Reporte egresados} es posible acceder directamente a secciones como análisis salarial, tecnologías más demandadas, satisfacción laboral, mapa de calor y métricas avanzadas; el sistema desplaza la vista hasta el gráfico seleccionado.

\begin{figure}[H]
\centering
\fbox{\parbox{0.85\textwidth}{\centering\vspace{2cm}[Insertar captura: dashboard administrador]\vspace{2cm}}}
\caption{Dashboard principal del administrador.}
\label{fig:dashboard-admin}
\end{figure}

\subsection{Gestión de usuarios egresados}

\subsubsection{Usuarios activos}

Ruta: \texttt{/admin/view-users}. La tabla lista los egresados registrados. Los elementos de la interfaz son:

\begin{enumerate}
\item \textbf{Buscador:} permite filtrar por nombre, correo o empresa.
\item \textbf{Botón limpiar filtros:} restablece la búsqueda.
\item \textbf{Botón actualizar:} recarga los datos desde el servidor.
\item \textbf{Tabla paginada:} muestra el detalle de cada egresado y acciones de edición o desactivación.
\end{enumerate}

\begin{figure}[H]
\centering
\fbox{\parbox{0.85\textwidth}{\centering\vspace{2cm}[Insertar captura: usuarios activos]\vspace{2cm}}}
\caption{Listado de usuarios egresados activos.}
\label{fig:usuarios-activos}
\end{figure}

\subsubsection{Usuarios eliminados}

Ruta: \texttt{/admin/view-deleted-users}. Muestra egresados desactivados y permite su reactivación.

\subsection{Gestión de administradores}

Rutas: \texttt{/admin/view-admin} (activos) y \texttt{/admin/view-deleted-admin} (eliminados). Permite administrar cuentas con privilegios administrativos de forma análoga a la gestión de egresados.

\subsection{Gestión de encuestas}

\subsubsection{Crear o editar encuesta}

Ruta: \texttt{/admin/create-encuesta}. El formulario se organiza en: (1) datos principales (título, estado, descripción, fechas); (2) definición de preguntas y tipos de respuesta; (3) guardado. Para editar una encuesta existente se utiliza \texttt{/admin/editar-encuesta/:id}.

\subsubsection{Listado y resultados}

\begin{itemize}
\item \textbf{Mis encuestas} (\texttt{/admin/view-encuesta}): listado y mantenimiento de encuestas.
\item \textbf{Resultados} (\texttt{/admin/view-encuestas-resultados}): análisis de respuestas, gráficos y exportación de datos.
\end{itemize}

\begin{figure}[H]
\centering
\fbox{\parbox{0.85\textwidth}{\centering\vspace{2cm}[Insertar captura: resultados encuesta]\vspace{2cm}}}
\caption{Vista de resultados de encuestas.}
\label{fig:resultados-encuesta}
\end{figure}

\subsection{Mi perfil del administrador}

Ruta: \texttt{/admin/mi-perfil}. Permite actualizar datos personales y credenciales del administrador.

\section{Módulos del sistema --- Usuario egresado}
\label{sec:manual-egresado}

\subsection{Completar perfil profesional}
\label{subsec:perfil-egresado}

Tras el primer ingreso, el sistema puede restringir el acceso al dashboard y a encuestas hasta completar el perfil (ruta \texttt{/user/mi-perfil}). Los campos obligatorios son: años de experiencia, nivel de educación, especialidad técnica, tipo de empleo actual, disponibilidad de cambio y área de interés.

\begin{figure}[H]
\centering
\fbox{\parbox{0.85\textwidth}{\centering\vspace{2cm}[Insertar captura: perfil incompleto]\vspace{2cm}}}
\caption{Solicitud de completar perfil profesional.}
\label{fig:perfil-incompleto}
\end{figure}

\subsection{Dashboard del egresado}

Ruta: \texttt{/user/dashboard}. Ofrece las mismas visualizaciones agregadas que el panel administrativo, sin la tarjeta de nuevos registros. El egresado puede navegar por las subsecciones del menú \emph{Reporte egresados}.

\subsection{Encuestas}

\subsubsection{Mis encuestas}

Ruta: \texttt{/user/view-encuestas}. Lista encuestas disponibles con buscador y filtro por estado. El botón \emph{Responder} conduce a \texttt{/user/responder-encuesta/:id}.

\subsubsection{Encuestas completadas}

Ruta: \texttt{/user/encuesta-completada}. Historial de encuestas ya enviadas.

\begin{figure}[H]
\centering
\fbox{\parbox{0.85\textwidth}{\centering\vspace{2cm}[Insertar captura: responder encuesta]\vspace{2cm}}}
\caption{Formulario de respuesta de encuesta.}
\label{fig:responder-encuesta}
\end{figure}

\subsection{Mi perfil del egresado}

Ruta: \texttt{/user/mi-perfil}. Permite editar y guardar la información profesional que alimenta los reportes del sistema.

\section{Resumen de rutas}

La Tabla~\ref{tab:rutas-manual} resume las rutas principales por rol.

\begin{table}[H]
\centering
\caption{Rutas principales del sistema SeguimientoApp.}
\label{tab:rutas-manual}
\begin{tabular}{|l|l|l|}
\hline
\textbf{Funcionalidad} & \textbf{Administrador} & \textbf{Egresado} \\
\hline
Inicio de sesión & \texttt{/auth} & \texttt{/auth} (LinkedIn) \\
Dashboard & \texttt{/admin/dashboard} & \texttt{/user/dashboard} \\
Mi perfil & \texttt{/admin/mi-perfil} & \texttt{/user/mi-perfil} \\
Usuarios & \texttt{/admin/view-users} & --- \\
Encuestas & \texttt{/admin/view-encuesta} & \texttt{/user/view-encuestas} \\
Resultados & \texttt{/admin/view-encuestas-resultados} & --- \\
Responder encuesta & --- & \texttt{/user/responder-encuesta/:id} \\
\hline
\end{tabular}
\end{table}

% Requiere en el preámbulo: \usepackage{float} \usepackage{graphicx} \usepackage{hyperref}

\chapter{Manual de usuarios}
\label{ap:manual-usuario}

\begin{center}
\textbf{MANUAL DE USO DEL SISTEMA DE SEGUIMIENTO DE EGRESADOS}\\[0.5em]
\textit{Trabajo de Título --- SeguimientoApp}
\end{center}

\section{Introducción}
\label{sec:manual-intro}

El objetivo de realizar este manual es poder ayudar a los diferentes usuarios que quieran utilizar este sistema, guiándolos en los pasos que debe realizar ya sea para ingresar al sistema como en el uso del mismo.

Este proyecto fue desarrollado con el fin de dar seguimiento a la trayectoria profesional de los egresados de la institución, consolidar información laboral y tecnológica, y apoyar la toma de decisiones académicas e institucionales mediante reportes y encuestas. El sistema contempla dos perfiles principales: \textbf{egresado} y \textbf{administrador}.

A continuación se detalla, en la primera sección, el ingreso al sistema; en las secciones siguientes, el funcionamiento de cada módulo según el rol del usuario.

\subsection{Ingreso al sistema}
\label{subsec:manual-ingreso}

\subsubsection{Requisitos y acceso general}

Para utilizar el sistema, el usuario debe abrir un navegador web y dirigirse a la URL de despliegue de la aplicación (por ejemplo, la ruta \texttt{/auth} del frontend). En la pantalla de inicio se podrá alternar entre el \textbf{Portal Egresados} y el \textbf{Panel Administrativo} mediante el botón de cambio de rol, como se muestra en la Figura~\ref{fig:login}.

\begin{figure}[H]
\centering
% \includegraphics[width=0.85\textwidth]{figuras/manual/login.png}
\fbox{\parbox{0.85\textwidth}{\centering\vspace{2cm}[Insertar captura: pantalla de login]\vspace{2cm}}}
\caption{Pantalla de inicio de sesión con selector de rol.}
\label{fig:login}
\end{figure}

\subsubsection{Ingreso del usuario egresado}

El acceso del egresado al sistema se realiza mediante autenticación con LinkedIn y contempla, para un usuario nuevo, un flujo de cuatro etapas que se describen a continuación: ingreso a la aplicación, autorización en LinkedIn, completitud del perfil profesional y acceso al dashboard principal.

\paragraph{Paso 1: Ingreso a la página del sistema.}

El primer paso consiste en abrir un navegador web y dirigirse a la URL de despliegue de la aplicación (por ejemplo, \texttt{http://localhost:4200} en ambiente de desarrollo, o la dirección institucional definida en producción). Al cargar el sitio, el sistema redirige automáticamente a la ruta de autenticación \texttt{/auth}, donde se presenta la pantalla denominada \emph{Portal Egresados}, como se muestra en la Figura~\ref{fig:login-egresado-paso1}.

En esta pantalla, el usuario debe verificar que \textbf{no} esté activa la vista de \emph{Panel Administrativo}. Si el formulario muestra campos de correo y contraseña, debe presionar el botón superior \emph{Egresados} para volver al modo correspondiente al perfil de egresado. Los elementos visibles en esta etapa son los siguientes:

\begin{enumerate}
\item \textbf{Encabezado del portal:} título \emph{Portal Egresados} y mensaje de bienvenida.
\item \textbf{Ícono de LinkedIn} junto al texto \emph{Conecta con tu perfil profesional}, que indica el método de acceso habilitado para este rol.
\item \textbf{Botón Ingresar con LinkedIn:} acción principal que inicia el proceso de autenticación externa.
\item \textbf{Botón de cambio de rol} (parte superior): permite alternar entre egresado y administrador; debe permanecer en modo egresado.
\end{enumerate}

\begin{figure}[H]
\centering
\fbox{\parbox{0.85\textwidth}{\centering\vspace{2cm}[Insertar captura: Portal Egresados --- Paso 1]\vspace{2cm}}}
\caption{Paso 1 --- Página de ingreso del egresado en la URL del sistema.}
\label{fig:login-egresado-paso1}
\end{figure}

\paragraph{Paso 2: Redirección y autenticación en LinkedIn.}

Al presionar el botón \emph{Ingresar con LinkedIn}, el sistema solicita al servidor la URL de autorización de LinkedIn y muestra un cuadro de diálogo con el mensaje \emph{Conectando con LinkedIn} o \emph{Redirigiendo a LinkedIn}, indicando que la operación está en curso. Posteriormente, el navegador abandona temporalmente la aplicación y carga el sitio oficial de LinkedIn (\texttt{linkedin.com}), donde el usuario deberá:

\begin{enumerate}
\item Iniciar sesión con su cuenta de LinkedIn, si aún no lo ha hecho en esa sesión del navegador.
\item Revisar los permisos que la aplicación solicita (acceso al perfil básico y datos profesionales).
\item Presionar el botón \emph{Permitir} o \emph{Aceptar} para autorizar la vinculación de su cuenta con el sistema.
\end{enumerate}

Una vez autorizado el acceso, LinkedIn redirige nuevamente a la aplicación con un código de autenticación. El sistema procesa dicho código en segundo plano, registra al usuario si es su primer ingreso o actualiza sus datos si ya existía previamente, y muestra un mensaje de éxito (\emph{Accediendo al portal...}). Finalmente, el egresado es conducido al área privada del sistema bajo la ruta \texttt{/user}, como se aprecia en la Figura~\ref{fig:login-egresado-paso2}.

\begin{figure}[H]
\centering
\fbox{\parbox{0.85\textwidth}{\centering\vspace{2cm}[Insertar captura: autorización LinkedIn --- Paso 2]\vspace{2cm}}}
\caption{Paso 2 --- Pantalla de autorización en LinkedIn y retorno al sistema.}
\label{fig:login-egresado-paso2}
\end{figure}

Si el usuario cancela la autorización o ocurre un error en la autenticación, el sistema permanece en la pantalla de login y despliega un mensaje de error descriptivo, permitiendo reintentar el proceso.

\paragraph{Paso 3: Completar el perfil profesional (usuario nuevo).}

Cuando el egresado ingresa por primera vez, el sistema importa desde LinkedIn información básica como nombre, fotografía, correo, empresa, cargo, ubicación y resumen profesional. Sin embargo, para generar métricas precisas del mercado laboral, la aplicación exige completar un conjunto de \textbf{campos obligatorios} antes de habilitar el dashboard y el módulo de encuestas.

Si el perfil no está completo, al intentar acceder a \texttt{/user/dashboard} o a las encuestas, el sistema redirige automáticamente a la ruta \texttt{/user/mi-perfil} y muestra un aviso destacado con el mensaje \emph{Complete su perfil profesional}, junto a una barra de progreso de completitud (ver Figura~\ref{fig:login-egresado-paso3}). Los campos que el usuario debe completar son:

\begin{enumerate}
\item Años de experiencia laboral.
\item Nivel de educación.
\item Especialidad técnica.
\item Tipo de empleo actual.
\item Disponibilidad para cambio de trabajo.
\item Área de interés.
\end{enumerate}

Además, el formulario incluye secciones de \emph{Información personal} (nombre, ubicación), \emph{Información profesional} (posición, empresa, industria, resumen) y \emph{Datos para métricas} (tecnologías, rango salarial, satisfacción laboral, entre otros campos opcionales). El usuario debe completar los campos marcados con asterisco rojo ($*$), presionar el botón \emph{Guardar cambios} y verificar que el indicador de progreso alcance el \textbf{100\%}, momento en el cual aparece el mensaje \emph{¡Perfil completado exitosamente!}

\begin{figure}[H]
\centering
\fbox{\parbox{0.85\textwidth}{\centering\vspace{2cm}[Insertar captura: completar perfil --- Paso 3]\vspace{2cm}}}
\caption{Paso 3 --- Completitud del perfil profesional para usuario nuevo.}
\label{fig:login-egresado-paso3}
\end{figure}

Hasta que el perfil no se encuentre completo, el menú lateral seguirá visible, pero las rutas protegidas del dashboard y encuestas permanecerán bloqueadas por el sistema de validación de perfil.

\paragraph{Paso 4: Acceso al Dashboard principal.}

Una vez completado el perfil, el egresado puede acceder al \emph{Dashboard principal} seleccionando en el menú lateral la opción \emph{Reporte egresados} $\rightarrow$ \emph{Dashboard principal}, o navegando directamente a \texttt{/user/dashboard}. En esta pantalla el usuario visualiza por primera vez el conjunto de indicadores agregados del mercado laboral tecnológico, alimentados con los datos consolidados de la comunidad de egresados (ver Figura~\ref{fig:login-egresado-paso4}).

Los elementos principales del dashboard son:

\begin{enumerate}
\item \textbf{Encabezado personalizado:} muestra la fotografía importada de LinkedIn, el nombre del egresado, su cargo actual, ubicación, empresa e industria. Al hacer clic en el nombre o la imagen, el usuario puede volver a \emph{Mi perfil}.
\item \textbf{Tarjetas de resumen del mercado:} indicadores como profesionales activos, satisfacción laboral promedio, tecnologías más utilizadas y empleabilidad del sector.
\item \textbf{Gráficos interactivos:} distribución salarial, tecnologías demandadas, análisis de experiencia, mapa de calor por industria y métricas avanzadas, accesibles también desde los submenús del panel lateral.
\item \textbf{Botón de menú (\texttt{☰}):} permite desplegar el panel de navegación para acceder a encuestas u otras secciones del reporte.
\item \textbf{Botón volver arriba:} aparece al desplazarse hacia abajo en páginas extensas de gráficos.
\end{enumerate}

\begin{figure}[H]
\centering
\fbox{\parbox{0.85\textwidth}{\centering\vspace{2cm}[Insertar captura: Dashboard principal --- Paso 4]\vspace{2cm}}}
\caption{Paso 4 --- Dashboard principal del egresado tras completar el perfil.}
\label{fig:login-egresado-paso4}
\end{figure}

A partir de este punto, el egresado dispone de acceso completo a las funcionalidades de su rol: consulta de reportes, respuesta de encuestas y actualización permanente de su perfil profesional (detalladas en la Sección~\ref{sec:manual-egresado}).

\subsubsection{Ingreso del usuario administrador}

El administrador debe activar la vista \emph{Panel Administrativo}, ingresar correo electrónico y contraseña, y presionar \emph{Iniciar Sesión}. La recuperación de contraseña está disponible mediante el enlace \emph{¿Olvidaste tu contraseña?}, que despliega un asistente en tres pasos: solicitud de código, verificación e ingreso de nueva contraseña.

\begin{figure}[H]
\centering
\fbox{\parbox{0.85\textwidth}{\centering\vspace{2cm}[Insertar captura: login administrador]\vspace{2cm}}}
\caption{Formulario de acceso administrativo.}
\label{fig:login-admin}
\end{figure}

\subsubsection{Navegación general}

Una vez autenticado, la interfaz se compone de una barra superior, un botón de menú (\texttt{☰}) que despliega el panel lateral, y el área de contenido. En el menú lateral, la tarjeta de perfil permite acceder a \emph{Mi perfil}. Las opciones del menú dependen del rol del usuario.

\begin{figure}[H]
\centering
\fbox{\parbox{0.85\textwidth}{\centering\vspace{2cm}[Insertar captura: layout principal]\vspace{2cm}}}
\caption{Vista general del sistema con menú lateral.}
\label{fig:layout}
\end{figure}

\section{Módulos del sistema --- Usuario administrador}
\label{sec:manual-admin}

\subsection{Dashboard y reporte de egresados}
\label{subsec:dashboard-admin}

Para ingresar al dashboard, el administrador debe seleccionar en el menú lateral la opción \emph{Reporte egresados} $\rightarrow$ \emph{Dashboard principal} (ruta \texttt{/admin/dashboard}).

La página presenta métricas agregadas del mercado laboral tecnológico. Los elementos principales son:

\begin{enumerate}
\item Encabezado con indicadores generales del mercado laboral.
\item Tarjeta de \emph{Nuevos registros} con la cantidad de egresados incorporados recientemente.
\item Conjunto de gráficos interactivos alimentados por los datos de los perfiles.
\end{enumerate}

Desde el submenú \emph{Reporte egresados} es posible acceder directamente a secciones como análisis salarial, tecnologías más demandadas, satisfacción laboral, mapa de calor y métricas avanzadas; el sistema desplaza la vista hasta el gráfico seleccionado.

\begin{figure}[H]
\centering
\fbox{\parbox{0.85\textwidth}{\centering\vspace{2cm}[Insertar captura: dashboard administrador]\vspace{2cm}}}
\caption{Dashboard principal del administrador.}
\label{fig:dashboard-admin}
\end{figure}

\subsection{Gestión de usuarios egresados}

\subsubsection{Usuarios activos}

Ruta: \texttt{/admin/view-users}. La tabla lista los egresados registrados. Los elementos de la interfaz son:

\begin{enumerate}
\item \textbf{Buscador:} permite filtrar por nombre, correo o empresa.
\item \textbf{Botón limpiar filtros:} restablece la búsqueda.
\item \textbf{Botón actualizar:} recarga los datos desde el servidor.
\item \textbf{Tabla paginada:} muestra el detalle de cada egresado y acciones de edición o desactivación.
\end{enumerate}

\begin{figure}[H]
\centering
\fbox{\parbox{0.85\textwidth}{\centering\vspace{2cm}[Insertar captura: usuarios activos]\vspace{2cm}}}
\caption{Listado de usuarios egresados activos.}
\label{fig:usuarios-activos}
\end{figure}

\subsubsection{Usuarios eliminados}

Ruta: \texttt{/admin/view-deleted-users}. Muestra egresados desactivados y permite su reactivación.

\subsection{Gestión de administradores}

Rutas: \texttt{/admin/view-admin} (activos) y \texttt{/admin/view-deleted-admin} (eliminados). Permite administrar cuentas con privilegios administrativos de forma análoga a la gestión de egresados.

\subsection{Gestión de encuestas}

\subsubsection{Crear o editar encuesta}

Ruta: \texttt{/admin/create-encuesta}. El formulario se organiza en: (1) datos principales (título, estado, descripción, fechas); (2) definición de preguntas y tipos de respuesta; (3) guardado. Para editar una encuesta existente se utiliza \texttt{/admin/editar-encuesta/:id}.

\subsubsection{Listado y resultados}

\begin{itemize}
\item \textbf{Mis encuestas} (\texttt{/admin/view-encuesta}): listado y mantenimiento de encuestas.
\item \textbf{Resultados} (\texttt{/admin/view-encuestas-resultados}): análisis de respuestas, gráficos y exportación de datos.
\end{itemize}

\begin{figure}[H]
\centering
\fbox{\parbox{0.85\textwidth}{\centering\vspace{2cm}[Insertar captura: resultados encuesta]\vspace{2cm}}}
\caption{Vista de resultados de encuestas.}
\label{fig:resultados-encuesta}
\end{figure}

\subsection{Mi perfil del administrador}

Ruta: \texttt{/admin/mi-perfil}. Permite actualizar datos personales y credenciales del administrador.

\section{Módulos del sistema --- Usuario egresado}
\label{sec:manual-egresado}

\subsection{Completar perfil profesional}
\label{subsec:perfil-egresado}

Tras el primer ingreso, el sistema puede restringir el acceso al dashboard y a encuestas hasta completar el perfil (ruta \texttt{/user/mi-perfil}). Los campos obligatorios son: años de experiencia, nivel de educación, especialidad técnica, tipo de empleo actual, disponibilidad de cambio y área de interés.

\begin{figure}[H]
\centering
\fbox{\parbox{0.85\textwidth}{\centering\vspace{2cm}[Insertar captura: perfil incompleto]\vspace{2cm}}}
\caption{Solicitud de completar perfil profesional.}
\label{fig:perfil-incompleto}
\end{figure}

\subsection{Dashboard del egresado}

Ruta: \texttt{/user/dashboard}. Ofrece las mismas visualizaciones agregadas que el panel administrativo, sin la tarjeta de nuevos registros. El egresado puede navegar por las subsecciones del menú \emph{Reporte egresados}.

\subsection{Encuestas}

\subsubsection{Mis encuestas}

Ruta: \texttt{/user/view-encuestas}. Lista encuestas disponibles con buscador y filtro por estado. El botón \emph{Responder} conduce a \texttt{/user/responder-encuesta/:id}.

\subsubsection{Encuestas completadas}

Ruta: \texttt{/user/encuesta-completada}. Historial de encuestas ya enviadas.

\begin{figure}[H]
\centering
\fbox{\parbox{0.85\textwidth}{\centering\vspace{2cm}[Insertar captura: responder encuesta]\vspace{2cm}}}
\caption{Formulario de respuesta de encuesta.}
\label{fig:responder-encuesta}
\end{figure}

\subsection{Mi perfil del egresado}

Ruta: \texttt{/user/mi-perfil}. Permite editar y guardar la información profesional que alimenta los reportes del sistema.

\section{Resumen de rutas}

La Tabla~\ref{tab:rutas-manual} resume las rutas principales por rol.

\begin{table}[H]
\centering
\caption{Rutas principales del sistema SeguimientoApp.}
\label{tab:rutas-manual}
\begin{tabular}{|l|l|l|}
\hline
\textbf{Funcionalidad} & \textbf{Administrador} & \textbf{Egresado} \\
\hline
Inicio de sesión & \texttt{/auth} & \texttt{/auth} (LinkedIn) \\
Dashboard & \texttt{/admin/dashboard} & \texttt{/user/dashboard} \\
Mi perfil & \texttt{/admin/mi-perfil} & \texttt{/user/mi-perfil} \\
Usuarios & \texttt{/admin/view-users} & --- \\
Encuestas & \texttt{/admin/view-encuesta} & \texttt{/user/view-encuestas} \\
Resultados & \texttt{/admin/view-encuestas-resultados} & --- \\
Responder encuesta & --- & \texttt{/user/responder-encuesta/:id} \\
\hline
\end{tabular}
\end{table}

% Requiere en el preámbulo: \usepackage{float} \usepackage{graphicx} \usepackage{hyperref}

\chapter{Manual de usuarios}
\label{ap:manual-usuario}

\begin{center}
\textbf{MANUAL DE USO DEL SISTEMA DE SEGUIMIENTO DE EGRESADOS}\\[0.5em]
\textit{Trabajo de Título --- SeguimientoApp}
\end{center}

\section{Introducción}
\label{sec:manual-intro}

El objetivo de realizar este manual es poder ayudar a los diferentes usuarios que quieran utilizar este sistema, guiándolos en los pasos que debe realizar ya sea para ingresar al sistema como en el uso del mismo.

Este proyecto fue desarrollado con el fin de dar seguimiento a la trayectoria profesional de los egresados de la institución, consolidar información laboral y tecnológica, y apoyar la toma de decisiones académicas e institucionales mediante reportes y encuestas. El sistema contempla dos perfiles principales: \textbf{egresado} y \textbf{administrador}.

A continuación se detalla, en la primera sección, el ingreso al sistema; en las secciones siguientes, el funcionamiento de cada módulo según el rol del usuario.

\subsection{Ingreso al sistema}
\label{subsec:manual-ingreso}

\subsubsection{Requisitos y acceso general}

Para utilizar el sistema, el usuario debe abrir un navegador web y dirigirse a la URL de despliegue de la aplicación (por ejemplo, la ruta \texttt{/auth} del frontend). En la pantalla de inicio se podrá alternar entre el \textbf{Portal Egresados} y el \textbf{Panel Administrativo} mediante el botón de cambio de rol, como se muestra en la Figura~\ref{fig:login}.

\begin{figure}[H]
\centering
% \includegraphics[width=0.85\textwidth]{figuras/manual/login.png}
\fbox{\parbox{0.85\textwidth}{\centering\vspace{2cm}[Insertar captura: pantalla de login]\vspace{2cm}}}
\caption{Pantalla de inicio de sesión con selector de rol.}
\label{fig:login}
\end{figure}

\subsubsection{Ingreso del usuario egresado}

El acceso del egresado al sistema se realiza mediante autenticación con LinkedIn y contempla, para un usuario nuevo, un flujo de cuatro etapas que se describen a continuación: ingreso a la aplicación, autorización en LinkedIn, completitud del perfil profesional y acceso al dashboard principal.

\paragraph{Paso 1: Ingreso a la página del sistema.}

El primer paso consiste en abrir un navegador web y dirigirse a la URL de despliegue de la aplicación (por ejemplo, \texttt{http://localhost:4200} en ambiente de desarrollo, o la dirección institucional definida en producción). Al cargar el sitio, el sistema redirige automáticamente a la ruta de autenticación \texttt{/auth}, donde se presenta la pantalla denominada \emph{Portal Egresados}, como se muestra en la Figura~\ref{fig:login-egresado-paso1}.

En esta pantalla, el usuario debe verificar que \textbf{no} esté activa la vista de \emph{Panel Administrativo}. Si el formulario muestra campos de correo y contraseña, debe presionar el botón superior \emph{Egresados} para volver al modo correspondiente al perfil de egresado. Los elementos visibles en esta etapa son los siguientes:

\begin{enumerate}
\item \textbf{Encabezado del portal:} título \emph{Portal Egresados} y mensaje de bienvenida.
\item \textbf{Ícono de LinkedIn} junto al texto \emph{Conecta con tu perfil profesional}, que indica el método de acceso habilitado para este rol.
\item \textbf{Botón Ingresar con LinkedIn:} acción principal que inicia el proceso de autenticación externa.
\item \textbf{Botón de cambio de rol} (parte superior): permite alternar entre egresado y administrador; debe permanecer en modo egresado.
\end{enumerate}

\begin{figure}[H]
\centering
\fbox{\parbox{0.85\textwidth}{\centering\vspace{2cm}[Insertar captura: Portal Egresados --- Paso 1]\vspace{2cm}}}
\caption{Paso 1 --- Página de ingreso del egresado en la URL del sistema.}
\label{fig:login-egresado-paso1}
\end{figure}

\paragraph{Paso 2: Redirección y autenticación en LinkedIn.}

Al presionar el botón \emph{Ingresar con LinkedIn}, el sistema solicita al servidor la URL de autorización de LinkedIn y muestra un cuadro de diálogo con el mensaje \emph{Conectando con LinkedIn} o \emph{Redirigiendo a LinkedIn}, indicando que la operación está en curso. Posteriormente, el navegador abandona temporalmente la aplicación y carga el sitio oficial de LinkedIn (\texttt{linkedin.com}), donde el usuario deberá:

\begin{enumerate}
\item Iniciar sesión con su cuenta de LinkedIn, si aún no lo ha hecho en esa sesión del navegador.
\item Revisar los permisos que la aplicación solicita (acceso al perfil básico y datos profesionales).
\item Presionar el botón \emph{Permitir} o \emph{Aceptar} para autorizar la vinculación de su cuenta con el sistema.
\end{enumerate}

Una vez autorizado el acceso, LinkedIn redirige nuevamente a la aplicación con un código de autenticación. El sistema procesa dicho código en segundo plano, registra al usuario si es su primer ingreso o actualiza sus datos si ya existía previamente, y muestra un mensaje de éxito (\emph{Accediendo al portal...}). Finalmente, el egresado es conducido al área privada del sistema bajo la ruta \texttt{/user}, como se aprecia en la Figura~\ref{fig:login-egresado-paso2}.

\begin{figure}[H]
\centering
\fbox{\parbox{0.85\textwidth}{\centering\vspace{2cm}[Insertar captura: autorización LinkedIn --- Paso 2]\vspace{2cm}}}
\caption{Paso 2 --- Pantalla de autorización en LinkedIn y retorno al sistema.}
\label{fig:login-egresado-paso2}
\end{figure}

Si el usuario cancela la autorización o ocurre un error en la autenticación, el sistema permanece en la pantalla de login y despliega un mensaje de error descriptivo, permitiendo reintentar el proceso.

\paragraph{Paso 3: Completar el perfil profesional (usuario nuevo).}

Cuando el egresado ingresa por primera vez, el sistema importa desde LinkedIn información básica como nombre, fotografía, correo, empresa, cargo, ubicación y resumen profesional. Sin embargo, para generar métricas precisas del mercado laboral, la aplicación exige completar un conjunto de \textbf{campos obligatorios} antes de habilitar el dashboard y el módulo de encuestas.

Si el perfil no está completo, al intentar acceder a \texttt{/user/dashboard} o a las encuestas, el sistema redirige automáticamente a la ruta \texttt{/user/mi-perfil} y muestra un aviso destacado con el mensaje \emph{Complete su perfil profesional}, junto a una barra de progreso de completitud (ver Figura~\ref{fig:login-egresado-paso3}). Los campos que el usuario debe completar son:

\begin{enumerate}
\item Años de experiencia laboral.
\item Nivel de educación.
\item Especialidad técnica.
\item Tipo de empleo actual.
\item Disponibilidad para cambio de trabajo.
\item Área de interés.
\end{enumerate}

Además, el formulario incluye secciones de \emph{Información personal} (nombre, ubicación), \emph{Información profesional} (posición, empresa, industria, resumen) y \emph{Datos para métricas} (tecnologías, rango salarial, satisfacción laboral, entre otros campos opcionales). El usuario debe completar los campos marcados con asterisco rojo ($*$), presionar el botón \emph{Guardar cambios} y verificar que el indicador de progreso alcance el \textbf{100\%}, momento en el cual aparece el mensaje \emph{¡Perfil completado exitosamente!}

\begin{figure}[H]
\centering
\fbox{\parbox{0.85\textwidth}{\centering\vspace{2cm}[Insertar captura: completar perfil --- Paso 3]\vspace{2cm}}}
\caption{Paso 3 --- Completitud del perfil profesional para usuario nuevo.}
\label{fig:login-egresado-paso3}
\end{figure}

Hasta que el perfil no se encuentre completo, el menú lateral seguirá visible, pero las rutas protegidas del dashboard y encuestas permanecerán bloqueadas por el sistema de validación de perfil.

\paragraph{Paso 4: Acceso al Dashboard principal.}

Una vez completado el perfil, el egresado puede acceder al \emph{Dashboard principal} seleccionando en el menú lateral la opción \emph{Reporte egresados} $\rightarrow$ \emph{Dashboard principal}, o navegando directamente a \texttt{/user/dashboard}. En esta pantalla el usuario visualiza por primera vez el conjunto de indicadores agregados del mercado laboral tecnológico, alimentados con los datos consolidados de la comunidad de egresados (ver Figura~\ref{fig:login-egresado-paso4}).

Los elementos principales del dashboard son:

\begin{enumerate}
\item \textbf{Encabezado personalizado:} muestra la fotografía importada de LinkedIn, el nombre del egresado, su cargo actual, ubicación, empresa e industria. Al hacer clic en el nombre o la imagen, el usuario puede volver a \emph{Mi perfil}.
\item \textbf{Tarjetas de resumen del mercado:} indicadores como profesionales activos, satisfacción laboral promedio, tecnologías más utilizadas y empleabilidad del sector.
\item \textbf{Gráficos interactivos:} distribución salarial, tecnologías demandadas, análisis de experiencia, mapa de calor por industria y métricas avanzadas, accesibles también desde los submenús del panel lateral.
\item \textbf{Botón de menú (\texttt{☰}):} permite desplegar el panel de navegación para acceder a encuestas u otras secciones del reporte.
\item \textbf{Botón volver arriba:} aparece al desplazarse hacia abajo en páginas extensas de gráficos.
\end{enumerate}

\begin{figure}[H]
\centering
\fbox{\parbox{0.85\textwidth}{\centering\vspace{2cm}[Insertar captura: Dashboard principal --- Paso 4]\vspace{2cm}}}
\caption{Paso 4 --- Dashboard principal del egresado tras completar el perfil.}
\label{fig:login-egresado-paso4}
\end{figure}

A partir de este punto, el egresado dispone de acceso completo a las funcionalidades de su rol: consulta de reportes, respuesta de encuestas y actualización permanente de su perfil profesional (detalladas en la Sección~\ref{sec:manual-egresado}).

\subsubsection{Ingreso del usuario administrador}

El administrador debe activar la vista \emph{Panel Administrativo}, ingresar correo electrónico y contraseña, y presionar \emph{Iniciar Sesión}. La recuperación de contraseña está disponible mediante el enlace \emph{¿Olvidaste tu contraseña?}, que despliega un asistente en tres pasos: solicitud de código, verificación e ingreso de nueva contraseña.

\begin{figure}[H]
\centering
\fbox{\parbox{0.85\textwidth}{\centering\vspace{2cm}[Insertar captura: login administrador]\vspace{2cm}}}
\caption{Formulario de acceso administrativo.}
\label{fig:login-admin}
\end{figure}

\subsubsection{Navegación general}

Una vez autenticado, la interfaz se compone de una barra superior, un botón de menú (\texttt{☰}) que despliega el panel lateral, y el área de contenido. En el menú lateral, la tarjeta de perfil permite acceder a \emph{Mi perfil}. Las opciones del menú dependen del rol del usuario.

\begin{figure}[H]
\centering
\fbox{\parbox{0.85\textwidth}{\centering\vspace{2cm}[Insertar captura: layout principal]\vspace{2cm}}}
\caption{Vista general del sistema con menú lateral.}
\label{fig:layout}
\end{figure}

\section{Módulos del sistema --- Usuario administrador}
\label{sec:manual-admin}

\subsection{Dashboard y reporte de egresados}
\label{subsec:dashboard-admin}

Para ingresar al dashboard, el administrador debe seleccionar en el menú lateral la opción \emph{Reporte egresados} $\rightarrow$ \emph{Dashboard principal} (ruta \texttt{/admin/dashboard}).

La página presenta métricas agregadas del mercado laboral tecnológico. Los elementos principales son:

\begin{enumerate}
\item Encabezado con indicadores generales del mercado laboral.
\item Tarjeta de \emph{Nuevos registros} con la cantidad de egresados incorporados recientemente.
\item Conjunto de gráficos interactivos alimentados por los datos de los perfiles.
\end{enumerate}

Desde el submenú \emph{Reporte egresados} es posible acceder directamente a secciones como análisis salarial, tecnologías más demandadas, satisfacción laboral, mapa de calor y métricas avanzadas; el sistema desplaza la vista hasta el gráfico seleccionado.

\begin{figure}[H]
\centering
\fbox{\parbox{0.85\textwidth}{\centering\vspace{2cm}[Insertar captura: dashboard administrador]\vspace{2cm}}}
\caption{Dashboard principal del administrador.}
\label{fig:dashboard-admin}
\end{figure}

\subsection{Gestión de usuarios egresados}

\subsubsection{Usuarios activos}

Ruta: \texttt{/admin/view-users}. La tabla lista los egresados registrados. Los elementos de la interfaz son:

\begin{enumerate}
\item \textbf{Buscador:} permite filtrar por nombre, correo o empresa.
\item \textbf{Botón limpiar filtros:} restablece la búsqueda.
\item \textbf{Botón actualizar:} recarga los datos desde el servidor.
\item \textbf{Tabla paginada:} muestra el detalle de cada egresado y acciones de edición o desactivación.
\end{enumerate}

\begin{figure}[H]
\centering
\fbox{\parbox{0.85\textwidth}{\centering\vspace{2cm}[Insertar captura: usuarios activos]\vspace{2cm}}}
\caption{Listado de usuarios egresados activos.}
\label{fig:usuarios-activos}
\end{figure}

\subsubsection{Usuarios eliminados}

Ruta: \texttt{/admin/view-deleted-users}. Muestra egresados desactivados y permite su reactivación.

\subsection{Gestión de administradores}

Rutas: \texttt{/admin/view-admin} (activos) y \texttt{/admin/view-deleted-admin} (eliminados). Permite administrar cuentas con privilegios administrativos de forma análoga a la gestión de egresados.

\subsection{Gestión de encuestas}

\subsubsection{Crear o editar encuesta}

Ruta: \texttt{/admin/create-encuesta}. El formulario se organiza en: (1) datos principales (título, estado, descripción, fechas); (2) definición de preguntas y tipos de respuesta; (3) guardado. Para editar una encuesta existente se utiliza \texttt{/admin/editar-encuesta/:id}.

\subsubsection{Listado y resultados}

\begin{itemize}
\item \textbf{Mis encuestas} (\texttt{/admin/view-encuesta}): listado y mantenimiento de encuestas.
\item \textbf{Resultados} (\texttt{/admin/view-encuestas-resultados}): análisis de respuestas, gráficos y exportación de datos.
\end{itemize}

\begin{figure}[H]
\centering
\fbox{\parbox{0.85\textwidth}{\centering\vspace{2cm}[Insertar captura: resultados encuesta]\vspace{2cm}}}
\caption{Vista de resultados de encuestas.}
\label{fig:resultados-encuesta}
\end{figure}

\subsection{Mi perfil del administrador}

Ruta: \texttt{/admin/mi-perfil}. Permite actualizar datos personales y credenciales del administrador.

\section{Módulos del sistema --- Usuario egresado}
\label{sec:manual-egresado}

\subsection{Completar perfil profesional}
\label{subsec:perfil-egresado}

Tras el primer ingreso, el sistema puede restringir el acceso al dashboard y a encuestas hasta completar el perfil (ruta \texttt{/user/mi-perfil}). Los campos obligatorios son: años de experiencia, nivel de educación, especialidad técnica, tipo de empleo actual, disponibilidad de cambio y área de interés.

\begin{figure}[H]
\centering
\fbox{\parbox{0.85\textwidth}{\centering\vspace{2cm}[Insertar captura: perfil incompleto]\vspace{2cm}}}
\caption{Solicitud de completar perfil profesional.}
\label{fig:perfil-incompleto}
\end{figure}

\subsection{Dashboard del egresado}

Ruta: \texttt{/user/dashboard}. Ofrece las mismas visualizaciones agregadas que el panel administrativo, sin la tarjeta de nuevos registros. El egresado puede navegar por las subsecciones del menú \emph{Reporte egresados}.

\subsection{Encuestas}

\subsubsection{Mis encuestas}

Ruta: \texttt{/user/view-encuestas}. Lista encuestas disponibles con buscador y filtro por estado. El botón \emph{Responder} conduce a \texttt{/user/responder-encuesta/:id}.

\subsubsection{Encuestas completadas}

Ruta: \texttt{/user/encuesta-completada}. Historial de encuestas ya enviadas.

\begin{figure}[H]
\centering
\fbox{\parbox{0.85\textwidth}{\centering\vspace{2cm}[Insertar captura: responder encuesta]\vspace{2cm}}}
\caption{Formulario de respuesta de encuesta.}
\label{fig:responder-encuesta}
\end{figure}

\subsection{Mi perfil del egresado}

Ruta: \texttt{/user/mi-perfil}. Permite editar y guardar la información profesional que alimenta los reportes del sistema.

\section{Resumen de rutas}

La Tabla~\ref{tab:rutas-manual} resume las rutas principales por rol.

\begin{table}[H]
\centering
\caption{Rutas principales del sistema SeguimientoApp.}
\label{tab:rutas-manual}
\begin{tabular}{|l|l|l|}
\hline
\textbf{Funcionalidad} & \textbf{Administrador} & \textbf{Egresado} \\
\hline
Inicio de sesión & \texttt{/auth} & \texttt{/auth} (LinkedIn) \\
Dashboard & \texttt{/admin/dashboard} & \texttt{/user/dashboard} \\
Mi perfil & \texttt{/admin/mi-perfil} & \texttt{/user/mi-perfil} \\
Usuarios & \texttt{/admin/view-users} & --- \\
Encuestas & \texttt{/admin/view-encuesta} & \texttt{/user/view-encuestas} \\
Resultados & \texttt{/admin/view-encuestas-resultados} & --- \\
Responder encuesta & --- & \texttt{/user/responder-encuesta/:id} \\
\hline
\end{tabular}
\end{table}
